\section{Held-out}

\begin{frame}[plain]
  \vfill
  \begin{center}
    {\color{brandaccent}\Large\textbf{HELD-OUT}}\\[0.4em]
    {\Large \NumNHeldout{} câu chưa từng bị chạm tới}\\[1em]
    {\small\color{brandmuted}
     Chạy \textbf{đúng một lần}, bằng cấu hình đã khóa từ trước.\\
     Đây --- và chỉ đây --- là con số được phép công bố.}
  \end{center}
  \vfill
\end{frame}

\begin{frame}{Vì sao phải có tập held-out}
  \small
  Mọi con số ở các slide trước đo trên \NumNDev{} câu development --- chính tập
  mà nhóm đã thử \textbf{4 prompt $\times$ 3 mức depth} rồi chọn cái điểm cao nhất.

  \begin{center}
    \colorbox{brandrule}{\parbox{0.88\textwidth}{\centering\vspace{0.3em}
      Chọn ra cái tốt nhất trong 12 lần thử, rồi báo cáo điểm của nó\\
      trên chính tập đã dùng để chọn, là \textbf{lệch lạc quan theo cấu trúc}.\vspace{0.3em}}}
  \end{center}

  \vspace{0.5em}
  \textbf{Quy trình đã tuân thủ, có dấu vết kiểm chứng được:}
  \begin{itemize}\small
    \item Tập \NumNHeldout{} câu / \NumNArticles{} bài chốt từ đầu bằng seed \NumHeldoutSeed{},
          băm lại khớp.
    \item Quyết định winner ký duyệt lúc \NumWinnerApprovedAt{}
          (\texttt{\NumHoDecisionSha\ldots}).
    \item Run held-out bắt đầu \NumHoStartedAt{} --- \textbf{sau} khi quyết định đã đóng băng.
    \item Coverage 100\%: \NumNHeldout{}/\NumNHeldout{} câu, 0 lỗi.
  \end{itemize}

  \takeaway{Thứ tự thời gian là bằng chứng: không thể chọn winner sau khi đã thấy điểm held-out.}
\end{frame}

\begin{frame}{Kết quả công bố --- \NumWinner{} trên \NumNHeldout{} câu held-out}
  \small
  \begin{tabularx}{\textwidth}{@{}L N N N@{}}
    \toprule
    & \textbf{Dev (\NumNDev)} & \textbf{Held-out (\NumNHeldout)} & \textbf{macro} \\
    \midrule
    Answer Correctness & \NumDevAC & \win{\NumHoAC} & \NumHoACMacro \\
    Exact Match & \NumDevEM & \pass{\NumHoEM} & --- \\
    Token F1 & \NumDevFone & \pass{\NumHoFone} & \NumHoFoneMacro \\
    Faithfulness & \NumDevFaith & \NumHoFaith & \NumHoFaithMacro \\
    Citation F1 & \NumDevCitFone & \NumHoCitFone & \NumHoCitFoneMacro \\
    Citation Validity & \NumDevCitVal & \NumHoCitVal & \NumHoCitValMacro \\
    \midrule
    \textbf{Hit@5} (truy xuất) & \NumDevHitFive & \fail{\NumHoHitFive} & --- \\
    \bottomrule
  \end{tabularx}

  \vspace{0.5em}
  Correctness chỉ tụt \textbf{0{,}0193}. Nhưng Hit@5 tụt \fail{\NumHoHitFiveDrop{}} ---
  tập held-out \textbf{khó hơn hẳn} ở tầng truy xuất.

  \vspace{0.3em}
  Micro và macro lệch nhau dưới 0{,}008 ở mọi metric $\rightarrow$ không bài báo nào chi phối.

  \takeaway{Chi phí toàn bộ run: \$\NumHoCost{}. Latency P50 \NumHoLatency{} ms.}
\end{frame}

\begin{frame}{Phát hiện quan trọng nhất: prompt \emph{không} bị overfit}
  \small
  Tách \NumNHeldout{} câu held-out theo việc truy xuất có đưa được bằng chứng ra hay không:

  \begin{tabularx}{\textwidth}{@{}L N N N@{}}
    \toprule
    & \textbf{n} & \textbf{Answer Corr.} & \textbf{Citation F1} \\
    \midrule
    Held-out, \texttt{gold\_in\_top5} & \NumHoNHit & \win{\NumHoHitAC} & \NumHoHitCitFone \\
    Held-out, \texttt{gold\_not\_in\_top5} & \NumHoNMiss & \fail{\NumHoMissAC} & \fail{\NumHoMissCitFone} \\
    \midrule
    \muted{Dev, \texttt{gold\_in\_top5}} & \muted{\NumNHit} & \muted{\NumPtwodFiveHitAC} & \muted{---} \\
    \bottomrule
  \end{tabularx}

  \vspace{0.6em}
  \begin{itemize}
    \item Trên các câu \textbf{truy xuất làm đúng việc}, held-out đạt \NumHoHitAC{}
          --- \textbf{cao hơn} dev (\NumPtwodFiveHitAC). Prompt tổng quát hóa được.
    \item Tỉ lệ truy xuất trượt tăng gần \textbf{ba lần}: \NumNMiss{}/\NumNDev{}
          $\rightarrow$ \NumHoNMiss{}/\NumNHeldout{}.
    \item Citation F1 trên nhóm trượt bằng \fail{\NumHoMissCitFone{}} --- theo định
          nghĩa: không có đoạn đúng nào trong context để mà trích dẫn.
  \end{itemize}

  \begin{center}
    \colorbox{brandrule}{\parbox{0.88\textwidth}{\centering\vspace{0.3em}
      \textbf{Toàn bộ khoảng cách dev $\rightarrow$ held-out là câu chuyện của truy xuất,\\
      không phải của prompt.}\vspace{0.3em}}}
  \end{center}

  \takeaway{Muốn cải thiện tiếp thì quay lại Phase 1, không phải viết lại prompt.}
\end{frame}
